\documentclass[12pt,a4paper]{article}

% --- 宏包加载 ---
\usepackage[utf8]{inputenc}
\usepackage[T1]{fontenc}
\usepackage{ctex} % 支持中文
\usepackage{amsmath, amssymb, amsfonts} % 数学公式
\usepackage{graphicx} % 插入图片
\usepackage{booktabs} % 表格样式
\usepackage{geometry} % 页面布局
\usepackage{cite} % 引用
\usepackage{hyperref} % 超链接
\usepackage{bm} % 数学加粗
\usepackage{listings} % 代码段
\usepackage{xcolor}

\geometry{left=2.5cm, right=2.5cm, top=3cm, bottom=3cm}

% --- 标题信息 ---
\title{\textbf{复杂双臂移动作业机器人的运动学建模与碰撞空间配置研究}}
\author{姓名：\underline{\hspace{2cm}} \quad 学号：\underline{\hspace{3cm}} \\ 专业：系统建模与仿真}
\date{\today}

\begin{document}

\maketitle

\begin{abstract}
随着智能制造和仓储物流的发展，具备高自由度的移动双臂机器人逐渐成为研究热点。本文针对一种集成升降旋转基座与双 6-DOF 机械臂的定制化复合机器人展开系统建模研究。首先，基于 URDF (Unified Robot Description Format) 对机器人进行树状拓扑结构的运动学建模。其次，针对高自由度系统带来的复杂自碰撞问题，基于 SRDF (Semantic Robot Description Format) 建立并优化了机器人的防碰撞矩阵 (ACM)。最后，结合 MoveIt 2 运动规划框架，构建了机器人的语义规划组与预定义状态。研究结果为后续的复杂空间避障与双臂协同轨迹规划奠定了系统模型基础。

\textbf{关键词：} 双臂移动机器人；运动学建模；碰撞检测；MoveIt 2；系统建模
\end{abstract}

\newpage

\section{引言}
在现代工业应用中，单臂机器人的工作空间往往受限于其安装基座。为了提高作业灵活性，移动双臂机器人（Mobile Dual-Arm Manipulator）通过集成移动底盘、升降机构与双臂，显著扩展了作业半径。然而，这类系统通常具有 10 个以上的自由度（DOF），属于典型的冗余自由度系统。建立精确的物理模型、运动学方程以及高效的碰撞避免机制是实现其自主作业的前提。

\section{机器人系统拓扑结构描述}
本研究对象由以下核心连杆 (Links) 与关节 (Joints) 组成，各组件通过父子关系构成树状运动链：

\subsection{基座运动系统}
基座系统包含底座 (\texttt{base\_link})、升降连杆 (\texttt{lifting\_link}) 与旋转连杆 (\texttt{revolute\_link})。
\begin{itemize}
    \item \textbf{升降关节 (\texttt{lifiting\_joint})}：棱柱关节 (Prismatic)，沿 Z 轴运动，行程 $d \in [0, 0.5]$ m。
    \item \textbf{旋转关节 (\texttt{revolute\_joint})}：转动关节 (Revolute)，绕 Z 轴旋转，范围 $\theta \in [-\pi, \pi]$。
\end{itemize}

\subsection{双臂子系统}
系统挂载两组 6-DOF 机械臂。以 Arm 1 为例，其相对于旋转基座的固定偏置为：
\begin{equation}
    P_{offset} = [0.0003, 0.0999, 0.3404]^T
\end{equation}
每个机械臂包含 6 个连续的旋转关节，末端安装有定制工具 (\texttt{tool.STL})。

\section{运动学系统建模}

\subsection{齐次变换矩阵定义}
根据 D-H 参数法或指数积公式，任意相邻连杆 $i-1$ 与 $i$ 之间的变换矩阵 $T_{i-1}^{i}$ 可表示为关节变量 $q_i$ 的函数：
\begin{equation}
T_{i-1}^{i}(q_i) = 
\begin{bmatrix}
\cos \theta_i & -\sin \theta_i & 0 & a_{i-1} \\
\sin \theta_i \cos \alpha_{i-1} & \cos \theta_i \cos \alpha_{i-1} & -\sin \alpha_{i-1} & -d_i \sin \alpha_{i-1} \\
\sin \theta_i \sin \alpha_{i-1} & \cos \theta_i \sin \alpha_{i-1} & \cos \alpha_{i-1} & d_i \cos \alpha_{i-1} \\
0 & 0 & 0 & 1
\end{bmatrix}
\end{equation}

\subsection{全系统正运动学}
对于 Arm 1 末端执行器相对于世界坐标系的位姿，其正运动学方程可表示为：
\begin{equation}
T_{world}^{tool1} = T_{world}^{base} \cdot T_{base}^{lift}(d_{lift}) \cdot T_{lift}^{rev}(\theta_{rev}) \cdot T_{rev}^{arm1\_base} \cdot \prod_{j=1}^{6} T_{j-1}^{j}(\theta_j) \cdot T_{6}^{tool}
\end{equation}

\section{碰撞模型优化策略}
在高自由度规划中，全量碰撞检测（Full Collision Check）的计算开销与连杆数量的平方成正比。

\subsection{允许碰撞矩阵 (ACM)}
基于 SRDF 的语义信息，我们排除了物理上不可能碰撞的连杆对。例如，\texttt{arm1\_base} 与 \texttt{revolute\_link} 属于相邻连杆，在模型层级将其设置为 \texttt{disable\_collisions}。
通过这种方式，碰撞检测对从 $N(N-1)/2$ 对大幅削减，有效降低了系统的实时计算负荷。

\section{MoveIt 规划系统配置}
为了实现任务层面的解耦，我们将 14-DOF 系统划分为以下规划组 (Planning Groups)：
\begin{table}[h]
\centering
\caption{机器人规划组定义}
\begin{tabular}{@{}lll@{}}
\toprule
规划组名称 & 自由度 & 包含核心关节 \\ \midrule
\texttt{arm1} & 6 & \texttt{arm1\_joint1} $\sim$ \texttt{arm1\_joint6} \\
\texttt{arm2} & 6 & \texttt{arm2\_joint1} $\sim$ \texttt{arm2\_joint6} \\
\texttt{agvBase} & 2 & \texttt{lifiting\_joint}, \texttt{revolute\_joint} \\ \bottomrule
\end{tabular}
\end{table}

\section{结论}
本文完成了对双臂移动作业机器人的物理结构分析与运动学参数建模。通过对自碰撞矩阵的裁剪优化，使得机器人在复杂的 ROS 2 环境下能够高效、安全地执行运动规划指令。该模型为后续的轨迹优化与双臂协同控制算法提供了可靠的仿真基础。

\end{document}
